	\begin{frame}
		\frametitle{DISCUSSION}
		\begin{itemize}
			\item In the multivariate Modified Poisson model, \textcolor {blue}{wealth index, type of place of residence, sex of the child, size of the child at birth, maternal education} were significantly associated with both Stunting and Underweight among children under five. 
			\item According to a study by Ambel in 2014, \textcolor {blue}{education determines the amount of health knowledge a mother has which is equally essential in child health}.  
			\item In urban areas, \textcolor {blue}{maternal health services such as ANC, PNC visits and health facility delivery are readily available and can equally raise community awareness and provide child immunization services and quality complementary feeding} (Bado \& Susuman, 2016) .
			\item Low income hoseholds have \textcolor {blue} {limited resources and thus less likely to access health facilities, expand their knowledge on access to health services and better nutritional practices that ensure children get a balanced diet} (Bhagowalia et al., 2012).  
		\end{itemize}
	\end{frame}
	\begin{frame}
		\frametitle{CONCLUSION}
		\begin{itemize}
			\item From the study, it was notable that the \textcolor {blue} {prevalence rates of malnutrition among under-fives were still a cause for concern and need to be addressed}. 
			\item Also, socio-economic  and bio-demographic characteristics such as type of place of residence, region, wealth-index , birth-order, maternal age, size of the child at birth and maternal education were found to be significant in determining a child’s nutrition status. 
			\item The study revealed that low income households, residing in the rural areas, with primary education level and lower contribute the highest cases of under-five malnutrition in Kenya. 
			\item Addressing these factors will likely reduce the cases of malnutrition. 
			\item \textcolor {blue} {The factors that influence malnutrition cases in Kenya are interrelated and thus there is no overall blanket program that can solve the problem at once}. 
			\item A more rigorous approach is looking at the problem from a multi-disciplinary perspective.
		\end{itemize}
	\end{frame}
	\begin{frame}
		\frametitle{RECOMMENDATIONS}
		\begin{itemize}
			\item A greater majority of the bio-demographic predictors that are maternal related are more personal. The most effective government intervention is \textcolor {blue}{civic education}. The Ministry of Health in conjunction with county governments should carry out campaigns , general advocacy and advertisements . 
			\item The region differences can be addressed by \textcolor {blue} {increasing the funding dedicated towards healthcare}. 
			\item There is need for further research that uses primary data where variables that were not captured are incorporated. 
		\end{itemize}
	\end{frame}
